\documentclass[conference]{IEEEtran}
\IEEEoverridecommandlockouts
% The preceding line is only needed to identify funding in the first footnote. If that is unneeded, please comment it out.
\usepackage{cite}
\usepackage{amsmath,amssymb,amsfonts}
\usepackage{algorithmic}
\usepackage{graphicx}
\usepackage{textcomp}
\usepackage{xcolor}
\def\BibTeX{{\rm B\kern-.05em{\sc i\kern-.025em b}\kern-.08em
    T\kern-.1667em\lower.7ex\hbox{E}\kern-.125emX}}
\begin{document}

\title{
\thanks{Identify applicable funding agency here. If none, delete this.}
}

\author{\IEEEauthorblockN{1\textsuperscript{st} Given Name Surname}
\IEEEauthorblockA{\textit{dept. name of organization (of Aff.)} \\
\textit{name of organization (of Aff.)}\\
City, Country \\
email address or ORCID}
\and
\IEEEauthorblockN{2\textsuperscript{nd} Given Name Surname}
\IEEEauthorblockA{\textit{dept. name of organization (of Aff.)} \\
\textit{name of organization (of Aff.)}\\
City, Country \\
email address or ORCID}}

\maketitle

\begin{abstract}
This survey provides a comprehensive overview of the current state of research in Multi-Camera Multi-Object Tracking (MCMOT) with trajectory and graph embedding. The field of MCMOT has garnered significant attention in recent years due to its wide range of applications in surveillance, security, and intelligent transportation systems. Our analysis of relevant papers reveals key findings, common themes, and contrasting viewpoints, highlighting the complexities and challenges associated with tracking multiple objects across multiple cameras.

The survey covers various aspects of MCMOT, including traditional methods, deep learning-based approaches, graph neural networks for 3D multi-object tracking, event-based frameworks, joint feature extractors, and ensemble training paradigms. We also discuss the state-of-the-art performance on benchmark datasets, challenges and limitations in multi-object tracking, and real-world applications of MCMOT.

The key themes that emerge from this survey include the importance of efficient data association across cameras and frames, the need for effective inter-camera tracking, and the potential of graph embedding techniques to improve MCMOT performance. Our analysis highlights the contributions of various approaches, including global graphical models, end-to-end methods, and improved similarity metrics.

This survey aims to provide a comprehensive understanding of the current state of research in MCMOT, highlighting both the challenges and limitations of existing algorithms and introducing novel approaches to address these issues. The scope of this survey is to summarize the key findings, contributions, and themes in MCMOT, with a focus on trajectory and graph embedding techniques. Our analysis provides insights into the future directions and open research challenges in this field, paving the way for further research and development in MCMOT.
\end{abstract}

\begin{IEEEkeywords}
large language models, multimodal learning, natural language processing, machine learning, artificial intelligence
\end{IEEEkeywords}


\section{Introduction to Multi-Object Tracking}
Introduction to Multi-Object Tracking

Multi-object tracking is a crucial aspect of computer vision, enabling the simultaneous tracking of multiple objects across one or more cameras. This field has garnered significant attention in recent years due to its applications in various domains, including surveillance, robotics, and autonomous vehicles. The challenge of multi-camera tracking is distinct from single-camera tracking, requiring new technologies and system architectures \cite{Challenge of Multi-Camera Tracking}. Researchers have proposed various approaches to address this challenge, including global graphical models, end-to-end methods, and graph convolutional networks.

Key findings in the field of multi-object tracking indicate that global graphical models can improve multi-camera object tracking performance by treating single-camera and inter-camera tracking similarities differently \cite{An equalised global graphical model-based approach for multi-camera object tracking}. End-to-end approaches, such as the Multi-Camera Tracking Transformer (MCTR), have also been shown to effectively track objects across multiple cameras with overlapping fields of view \cite{MCTR: Multi Camera Tracking Transformer}. Furthermore, pipelines that model tracking problems as global graphs and adopt optimization algorithms, such as Generalized Maximum Multi Clique, can efficiently associate data across cameras and frames \cite{Multi-camera Multi-Object Tracking}. The use of graph convolutional networks (GCNs) has also been explored for associating single-camera trajectories into multi-camera global trajectories, providing a robust solution to large camera unsynchronizations \cite{Graph Convolutional Network for Multi-Target Multi-Camera Vehicle Tracking}.

The need for efficient data association across cameras and frames is a common challenge in multi-camera tracking \cite{Challenge of Multi-Camera Tracking, Multi-camera Multi-Object Tracking}. Researchers have employed various techniques to address this challenge, including the use of appearance and motion features, such as LOMO and Hankel matrix \cite{Multi-camera Multi-Object Tracking}. Global graphical models and optimization algorithms are popular approaches to addressing the data association challenge \cite{An equalised global graphical model-based approach for multi-camera object tracking, Multi-camera Multi-Object Tracking}. End-to-end methods, such as MCTR, are also gaining interest in multi-camera tracking \cite{MCTR: Multi Camera Tracking Transformer}.

Contrasting viewpoints or approaches exist in the field of multi-object tracking. For instance, heuristic techniques are still widely used, while end-to-end approaches like MCTR are gaining popularity \cite{MCTR: Multi Camera Tracking Transformer}. Additionally, most tracking methods focus on single-camera tracking, but multi-camera tracking requires different approaches and techniques \cite{Challenge of Multi-Camera Tracking, An equalised global graphical model-based approach for multi-camera object tracking}. Different optimization algorithms, such as Generalized Maximum Multi Clique and GCNs, are being used to address the challenge of data association in multi-camera tracking \cite{Multi-camera Multi-Object Tracking, Graph Convolutional Network for Multi-Target Multi-Camera Vehicle Tracking}.

From a technical standpoint, global graphical models treat single-camera and inter-camera tracking similarities differently and use optimization algorithms to associate data across cameras and frames \cite{An equalised global graphical model-based approach for multi-camera object tracking, Multi-camera Multi-Object Tracking}. End-to-end methods, such as MCTR, use end-to-end detectors like DETR to produce detections and detection embeddings independently for each camera view \cite{MCTR: Multi Camera Tracking Transformer}. GCNs are used to associate single-camera trajectories into multi-camera global trajectories, providing a robust solution to large camera unsynchronizations \cite{Graph Convolutional Network for Multi-Target Multi-Camera Vehicle Tracking}. Appearance and motion features are extracted using techniques such as LOMO and Hankel matrix \cite{Multi-camera Multi-Object Tracking}.

Despite the advancements in multi-object tracking, several limitations and challenges remain. Efficient data association across cameras and frames continues to be a significant challenge \cite{Challenge of Multi-Camera Tracking, Multi-camera Multi-Object Tracking}. Scalability is also a concern, as multi-camera tracking systems can be computationally expensive and may require large amounts of training data \cite{MCTR: Multi Camera Tracking Transformer}. Large camera unsynchronizations can affect the performance of multi-camera tracking systems \cite{Graph Convolutional Network for Multi-Target Multi-Camera Vehicle Tracking}. Class imbalance can occur in multi-target multi-camera tracking, requiring specialized loss functions to address \cite{Graph Convolutional Network for Multi-Target Multi-Camera Vehicle Tracking}. Addressing these challenges will be crucial to the development of robust and efficient multi-object tracking systems.

\section{Background and Motivation}
The problem of Multi-Camera Multi-Object Tracking (MCMOT) has garnered significant attention in recent years due to its wide range of applications in surveillance, security, and intelligent transportation systems. A thorough analysis of relevant papers reveals key findings, common themes, contrasting viewpoints, technical details, and limitations that have shaped the current state of MCMOT research. For instance, the paper "Challenge of Multi-Camera Tracking" \cite{Challenge of Multi-Camera Tracking} summarizes the problems in multi-camera tracking and identifies new directions for future work, highlighting the need for efficient data association across cameras and frames.

The use of graph-based models to represent the relationships between objects across cameras is a common approach, as seen in "An equalised global graphical model-based approach for multi-camera object tracking" \cite{An equalised global graphical model-based approach for multi-camera object tracking} and "Multi-camera Multi-Object Tracking" \cite{Multi-camera Multi-Object Tracking}. These papers propose global graph models with improved similarity metrics to enhance overall multi-camera object tracking performance. Furthermore, the introduction of end-to-end approaches, such as "MCTR: Multi Camera Tracking Transformer" \cite{MCTR: Multi Camera Tracking Transformer}, has tailored solutions for multi-object detection and tracking across multiple cameras with overlapping fields of view.

A closer examination of these papers reveals that the need for efficient data association across cameras and frames is a recurring theme in MCMOT research. The use of appearance and dynamic motion similarities to compute good similarity scores for tracking is also a common approach \cite{Graph Convolutional Network for Multi-Target Multi-Camera Vehicle Tracking}. However, dealing with class imbalance and large camera unsynchronizations remains a significant challenge, as highlighted in several papers \cite{Challenge of Multi-Camera Tracking, An equalised global graphical model-based approach for multi-camera object tracking}.

The development of robust and scalable algorithms that can handle real-world scenarios with multiple objects, cameras, and occlusions is an ongoing challenge. Various feature extraction algorithms, such as Local Maximal Occurrence Representation (LOMO) and Hankel matrix, have been employed to address this issue \cite{Multi-camera Multi-Object Tracking}. Additionally, different optimization problems, including Generalized Maximum Multi Clique optimization problem and Iterative Hankel Total Least Squares (IHTLS), have been used to solve the MCMOT problem \cite{An equalised global graphical model-based approach for multi-camera object tracking, Graph Convolutional Network for Multi-Target Multi-Camera Vehicle Tracking}.

The trend towards using deep learning-based methods in MCMOT research is also evident, with the use of end-to-end detectors like DEtector TRansformer (DETR) and Graph Convolutional Network \cite{MCTR: Multi Camera Tracking Transformer, Graph Convolutional Network for Multi-Target Multi-Camera Vehicle Tracking}. These approaches have shown promise in improving the efficiency and accuracy of MCMOT systems. However, further research is needed to develop more efficient, robust, and scalable algorithms that can handle real-world scenarios.

In conclusion, the analysis of these papers highlights the complexity and challenges of the MCMOT problem. While significant progress has been made in recent years, there is still a need for further research to address the limitations and challenges associated with MCMOT. The use of graph-based models, end-to-end methods, and deep learning-based approaches reflects the trend towards using more advanced techniques to solve the MCMOT problem, and their implications and connections have the potential to significantly impact the development of intelligent surveillance systems.

\section{Taxonomy of Multi-Object Tracking Approaches}
The taxonomy of multi-object tracking approaches is a crucial aspect of understanding the complexities and challenges associated with tracking multiple objects across multiple cameras. Based on the analysis of relevant papers, it can be observed that the field has evolved significantly over the years, with various approaches being proposed to address the challenges of camera unsynchronizations, object re-identification, and large-scale tracking \cite{An equalised global graphical model-based approach for multi-camera object tracking}. One key theme that emerges from the literature is the importance of global solutions that consider both single-camera and inter-camera tracking to achieve robust multi-camera object tracking performance \cite{Multi-Camera Tracking Transformer}. 

The use of graphical models, such as global graphical models, has been shown to improve multi-camera object tracking performance by considering both single-camera and inter-camera tracking \cite{An equalised global graphical model-based approach for multi-camera object tracking}. End-to-end approaches, such as the Multi-Camera Tracking Transformer (MCTR), have also been proposed as effective solutions for multi-object detection and tracking across multiple cameras with overlapping fields of view \cite{Multi-Camera Tracking Transformer}. Furthermore, Graph Convolutional Networks have been used to associate single-camera trajectories into multi-camera global trajectories, providing a robust solution to large camera unsynchronizations \cite{Graph Convolutional Network for Multi-Target Multi-Camera Vehicle Tracking}.

A tracked target-based taxonomy can be used to present object tracking algorithms, highlighting advantages and limitations of different approaches \cite{Object Tracking in Videos: Approaches and Issues}. The literature also highlights the challenges of multi-camera tracking, including camera unsynchronizations, varying fields of view, and object re-identification. Most papers emphasize the need for global solutions that consider both single-camera and inter-camera tracking to achieve robust multi-camera object tracking performance \cite{An equalised global graphical model-based approach for multi-camera object tracking}. The importance of graph-based methods, such as graphical models and Graph Convolutional Networks, is also evident in the literature, as they are commonly used to model relationships between objects across cameras \cite{Graph Convolutional Network for Multi-Target Multi-Camera Vehicle Tracking}.

The analysis of papers also reveals contrasting viewpoints or approaches, including heuristic vs. end-to-end methods and single-camera vs. global tracking \cite{Multi-Camera Tracking Transformer}. While some papers focus on single-camera object tracking, others emphasize the need for global solutions that consider multiple cameras \cite{An equalised global graphical model-based approach for multi-camera object tracking}. The technical details and methodologies employed in the literature include the use of graphical models, transformer architectures, and Graph Convolutional Networks \cite{Graph Convolutional Network for Multi-Target Multi-Camera Vehicle Tracking}. 

The limitations and challenges highlighted in the literature include camera unsynchronizations, class imbalance, and the need for large-scale datasets to train and evaluate models \cite{Graph Convolutional Network for Multi-Target Multi-Camera Vehicle Tracking}. These findings have significant implications for the development of robust multi-camera object tracking algorithms, emphasizing the need for global solutions that consider both single-camera and inter-camera tracking. The connections between different approaches and techniques also highlight the importance of continued research in this area to address the challenges and limitations associated with multi-camera multi-object tracking \cite{Object Tracking in Videos: Approaches and Issues}. Overall, the taxonomy of multi-object tracking approaches provides a comprehensive overview of the current state-of-the-art in multi-camera multi-object tracking, highlighting key themes, developments, and implications for future research.

\section{Traditional Methods for Multi-Object Tracking}
Traditional methods for multi-object tracking have been extensively explored in the literature, with a focus on addressing the challenges posed by multi-camera systems. The development of these traditional methods has led to significant advancements in the field, with key findings including the proposal of global graph models to improve multi-camera object tracking performance \cite{Title of paper 2} \cite{Title of paper 5}. These models treat single camera tracking and inter-camera tracking differently, allowing for more accurate and robust tracking. Additionally, end-to-end approaches have been introduced, such as MCTR, which leverages detectors like DETR to produce detections and detection embeddings independently for each camera view \cite{Title of paper 3}.

The use of graph models has also been a common theme in traditional methods for multi-object tracking. For example, Graph Convolutional Networks have been proposed to associate single-camera trajectories into multi-camera global trajectories, providing a global solution robust to large camera unsynchronizations \cite{Title of paper 4}. Furthermore, the importance of efficient data association across cameras and frames has been emphasized in several papers, including \cite{Title of paper 1} \cite{Title of paper 2} \cite{Title of paper 5}. This highlights the need for effective methods to associate objects across different camera views and time steps.

The traditional methods for multi-object tracking also exhibit contrasting viewpoints or approaches. For instance, papers \cite{Title of paper 2} and \cite{Title of paper 5} employ traditional graph-based methods, while paper \cite{Title of paper 3} introduces an end-to-end approach using detectors like DETR. Moreover, papers \cite{Title of paper 2} and \cite{Title of paper 4} focus on global optimization using graph models, whereas paper \cite{Title of paper 3} updates track embeddings locally at each frame. The choice of features also varies, with paper \cite{Title of paper 5} using hand-crafted features (LOMO) for appearance representation, while papers \cite{Title of paper 3} and \cite{Title of paper 4} learn representations using detectors like DETR or Graph Convolutional Networks.

The technical details and methodologies employed in traditional methods for multi-object tracking include graph construction, similarity metrics, and optimization techniques. For example, papers \cite{Title of paper 2} \cite{Title of paper 4} \cite{Title of paper 5} construct graphs to represent relationships between objects and cameras, while papers \cite{Title of paper 2} and \cite{Title of paper 5} propose improved similarity metrics for inter-camera object tracking. Additionally, papers \cite{Title of paper 2} \cite{Title of paper 3} \cite{Title of paper 4} employ optimization techniques, such as global graph models, end-to-end detectors, or Graph Convolutional Networks, to achieve accurate multi-object tracking.

Despite the advancements in traditional methods for multi-object tracking, several limitations and challenges remain. For instance, multi-camera systems with large numbers of cameras and objects pose scalability challenges. Large camera unsynchronizations can also significantly impact tracking performance \cite{Title of paper 4}. Furthermore, class imbalance issues can arise in multi-target multi-camera vehicle tracking, which can be addressed using custom loss functions \cite{Title of paper 4}. The need for large-scale annotated datasets to train and evaluate these methods is also a significant limitation, as it can be time-consuming and expensive to obtain \cite{Title of paper 3} \cite{Title of paper 5}. Overall, traditional methods for multi-object tracking have made significant progress in addressing the challenges posed by multi-camera systems, but further research is needed to overcome the remaining limitations and challenges.

\section{Deep Learning-Based Methods for Multi-Object Tracking}
Deep learning-based methods have revolutionized the field of multi-object tracking, offering significant improvements in accuracy and efficiency. Recent advancements in this area have been profound, with various techniques being proposed to address the complexities of tracking multiple objects across multiple cameras. This section provides an overview of the key findings and contributions of deep learning-based methods for multi-object tracking, highlighting common themes and patterns, contrasting viewpoints or approaches, technical details and methodologies, limitations, and challenges.

The introduction of the Multi-Camera Tracking Transformer (MCTR) \cite{Title of MCTR paper} has been a significant development in this field, offering an end-to-end approach for multi-camera multi-object tracking. This method leverages transformers to produce detections and detection embeddings independently for each camera view, demonstrating the potential of deep learning-based architectures in handling complex tracking scenarios. Another notable contribution is the Graph Convolutional Network for Multi-Target Multi-Camera Vehicle Tracking \cite{Title of Graph Convolutional Network paper}, which proposes a graph convolutional network to associate single-camera trajectories into multi-camera global trajectories. This approach has shown robustness to large camera unsynchronizations, highlighting the importance of graph-based models in multi-object tracking.

A comprehensive survey of embedding methods for multi-object tracking \cite{Title of Recent Advances in Embedding Methods paper} provides valuable insights into the various techniques employed in this field. The authors analyze and summarize existing approaches from seven different perspectives, emphasizing the diversity of embedding methods used in multi-object tracking. Furthermore, an equalised global graphical model-based approach for multi-camera object tracking \cite{Title of An equalised global graphical model-based approach paper} improves overall multi-camera object tracking performance by treating single camera tracking and inter-camera tracking similarities differently in a global graph model.

A common theme among these papers is the reliance on deep learning-based approaches, which have become the cornerstone of modern multi-object tracking systems. The use of multi-camera tracking is also prevalent, with many papers addressing the challenges of handling multiple cameras with overlapping or non-overlapping fields of view. Embedding methods are a crucial component of many multi-object tracking approaches, with various papers exploring different types of embeddings, such as patch-level, single-frame, and cross-frame joint embedding. Graph-based models are also widely used to represent relationships between objects and cameras, demonstrating their effectiveness in multi-object tracking.

However, there are contrasting viewpoints or approaches among these papers. For instance, MCTR proposes an end-to-end approach for multi-camera multi-object tracking, while other papers, such as the Graph Convolutional Network, adopt a two-stage approach, first detecting objects and then associating them across cameras. The choice of embedding strategy also varies, with different papers employing distinct methods, such as patch-level, single-frame, or cross-frame joint embedding. Additionally, some papers focus on global optimization, while others prioritize local optimization methods.

From a technical standpoint, the papers employ various architectures, including transformer architectures \cite{Title of MCTR paper} and graph convolutional networks \cite{Title of Graph Convolutional Network paper}. Custom loss functions and optimization methods are also proposed to address specific challenges in multi-object tracking. For example, MCTR's soft probabilistic association and the Graph Convolutional Network's new loss function demonstrate the importance of tailored solutions for this complex task.

Despite the significant advancements in deep learning-based methods for multi-object tracking, several limitations and challenges remain. Scalability and efficiency are major concerns, as multi-camera multi-object tracking can be computationally expensive and challenging to scale to large camera networks. Class imbalance and occlusions are also significant challenges, with Graph Convolutional Network for Multi-Target Multi-Camera Vehicle Tracking addressing class imbalance and other papers discussing the difficulties posed by occlusions and appearance changes. The need for large-scale datasets is also emphasized, as many papers highlight the importance of extensive training and evaluation data to develop effective multi-object tracking methods.

In conclusion, deep learning-based methods have transformed the field of multi-object tracking, offering improved accuracy and efficiency in complex tracking scenarios. By analyzing the key findings and contributions of recent papers, we can identify common themes and patterns, contrasting viewpoints or approaches, technical details and methodologies, limitations, and challenges. The development of more effective and efficient approaches to multi-object tracking will rely on addressing these challenges and leveraging the strengths of deep learning-based architectures, ultimately informing the creation of robust and scalable tracking systems for real-world applications.

\section{Graph Neural Networks for 3D Multi-Object Tracking}
Graph Neural Networks for 3D Multi-Object Tracking have emerged as a promising approach in recent years, with various architectures being proposed to tackle the challenges of tracking multiple objects in 3D space. The use of Graph Convolutional Network (GCN), global graphical models, and time domain GCN has been explored in several papers \cite{Title of paper 1}, \cite{Title of paper 2}, and \cite{Title of paper 5}. These contributions have improved the overall performance of multi-camera object tracking by leveraging the strengths of graph neural networks in modeling complex relationships between objects.

A key finding in this area is the importance of attention mechanisms in improving tracking performance, as highlighted in papers such as \cite{Title of paper 3} and \cite{Title of paper 4}. These works propose attention-based approaches to mitigate modality intermittence in 3D multi-object tracking and to focus on relevant objects. The use of spatio-temporal representations is another common theme, with papers proposing frameworks that integrate spatial and temporal information from multiple sensors \cite{Title of paper 4}.

The technical details of the proposed GNN architectures vary, with some using GCN, global graphical models, or multi-modal GNNs. The methodologies used to evaluate the performance of these approaches include using benchmark datasets such as nuScenes and KITTI, as well as evaluating metrics such as AMOTA, MOTA, and detection accuracy \cite{Title of paper 1}, \cite{Title of paper 3}, \cite{Title of paper 4}, and \cite{Title of paper 5}. While some papers focus on developing global graph models for multi-camera object tracking \cite{Title of paper 2}, others propose GNN architectures that can handle multiple modalities and mitigate modality intermittence \cite{Title of paper 3}.

Despite the promising results, there are limitations and challenges associated with these approaches. The need for large-scale labeled datasets to train GNNs is a significant limitation, as highlighted in \cite{Title of paper 3}. Additionally, handling class imbalance and modality intermittence in multi-object tracking remains a challenge \cite{Title of paper 1} and \cite{Title of paper 3}. The computational complexity of some proposed approaches may also limit their applicability in real-time systems \cite{Title of paper 4}.

Future research in this area should focus on developing more efficient and effective GNN architectures for multi-object tracking, improving the robustness of GNNs to various types of noise and interference, and exploring the application of GNNs to other related tasks such as object detection and segmentation. The connections between graph neural networks and other machine learning approaches, such as deep learning and reinforcement learning, should also be further explored to leverage their strengths in tackling complex tracking problems. Overall, the use of Graph Neural Networks for 3D Multi-Object Tracking has shown significant promise, and continued research in this area is expected to lead to even more effective and efficient tracking systems.

\section{Event-Based Frameworks for Dynamic Object Tracking}
Event-Based Frameworks for Dynamic Object Tracking have garnered significant attention in recent years, owing to their potential to efficiently track objects across multiple cameras while leveraging trajectory and graph embedding techniques \cite{MCTR: Multi-Camera Tracking Transformer}. This paradigm shift is largely attributed to the advent of event-based cameras that can capture high-dynamic range images with low latency, thereby facilitating real-time object tracking \cite{e-TLD: Event-based Framework for Dynamic Object Tracking}. The papers analyzed in this section present several key findings and contributions to the field, including the introduction of end-to-end approaches for multi-camera tracking \cite{MCTR: Multi-Camera Tracking Transformer}, long-term object tracking frameworks using event cameras \cite{e-TLD: Event-based Framework for Dynamic Object Tracking}, and pipelines for multi-target visual tracking under multi-camera systems \cite{Multi-camera Multi-Object Tracking}.

A common theme that emerges across these papers is the use of graph models to represent the tracking problem, highlighting the importance of this approach in multi-camera multi-object tracking \cite{An equalised global graphical model-based approach for multi-camera object tracking}. Real-time processing is another crucial aspect, with techniques such as Bayesian bootstrapping and direct image alignment being employed to achieve efficient tracking \cite{Tracking 6-DoF Object Motion from Events and Frames}. The combination of appearance and motion cues is also a prevalent strategy, using techniques such as LOMO feature extraction and Hankel matrix construction to improve tracking performance \cite{Multi-camera Multi-Object Tracking}.

However, contrasting viewpoints or approaches are also evident, with some papers advocating for end-to-end approaches \cite{MCTR: Multi-Camera Tracking Transformer}, while others propose pipeline-based approaches \cite{Multi-camera Multi-Object Tracking}. The choice between single camera and multi-camera tracking is another point of divergence, with some papers focusing on the former \cite{e-TLD: Event-based Framework for Dynamic Object Tracking} and others emphasizing the importance of inter-camera tracking \cite{An equalised global graphical model-based approach for multi-camera object tracking}. Furthermore, the use of event-based versus frame-based cameras is a topic of debate, with some papers combining both types of cameras \cite{Tracking 6-DoF Object Motion from Events and Frames}.

From a technical standpoint, the papers employ a range of methodologies, including transformer architectures \cite{MCTR: Multi-Camera Tracking Transformer}, online learning \cite{e-TLD: Event-based Framework for Dynamic Object Tracking}, and Generalized Maximum Multi Clique optimization \cite{Multi-camera Multi-Object Tracking}. Despite these advances, several limitations and challenges remain, such as scalability, occlusion and clutter, real-time processing, and the need for standardized evaluation metrics \cite{MCTR: Multi-Camera Tracking Transformer}. These findings have significant implications for the development of event-based frameworks for dynamic object tracking, highlighting the need for further research into robust and efficient tracking algorithms that can operate in complex scenarios.

The connections between these papers are evident, with each contribution building upon or complementing existing work in the field. For instance, the use of graph models in \cite{An equalised global graphical model-based approach for multi-camera object tracking} is closely related to the pipeline proposed in \cite{Multi-camera Multi-Object Tracking}. Similarly, the event-based framework presented in \cite{e-TLD: Event-based Framework for Dynamic Object Tracking} shares similarities with the approach outlined in \cite{Tracking 6-DoF Object Motion from Events and Frames}. These connections underscore the importance of continued research into event-based frameworks for dynamic object tracking, with a focus on addressing the limitations and challenges identified in these papers. Ultimately, the development of robust and efficient tracking algorithms will have significant implications for a range of applications, including surveillance, robotics, and autonomous vehicles \cite{MCTR: Multi-Camera Tracking Transformer}.

\section{Joint Feature Extractors and Ensemble Training Paradigms}
The development of joint feature extractors and ensemble training paradigms has been a crucial aspect of advancing multi-camera multi-object tracking (MCMOT) systems. Recent studies have proposed various approaches to improve the accuracy and efficiency of MCMOT, including graph convolutional networks \cite{Graph Convolutional Network for Multi-Target Multi-Camera Vehicle Tracking}, global graph models \cite{Multi-camera Multi-Object Tracking}, and end-to-end transformer-based methods \cite{MCTR: Multi Camera Tracking Transformer}. These approaches have demonstrated significant improvements in handling complex scenarios, such as camera unsynchronizations, class imbalance, and efficient data association across cameras.

A key theme emerging from these studies is the utilization of graph-based models to represent relationships between objects across cameras. For instance, the graph convolutional network proposed in \cite{Graph Convolutional Network for Multi-Target Multi-Camera Vehicle Tracking} associates single-camera trajectories into multi-camera global trajectories, outperforming related work. Similarly, the global graph model introduced in \cite{Multi-camera Multi-Object Tracking} takes into account both across-frame and across-camera data correlation using Generalized Maximum Multi Clique optimization. The use of graph-based models enables the effective capture of complex interactions between objects and cameras, leading to improved tracking performance.

Another significant development is the proposal of end-to-end approaches for MCMOT, such as the transformer-based method presented in \cite{MCTR: Multi Camera Tracking Transformer}. This approach simplifies the tracking process by jointly optimizing detection and tracking tasks, resulting in improved performance and efficiency. In contrast, heuristic methods, which rely on hand-designed rules and thresholds, have been criticized for their limitations in handling complex scenarios \cite{Challenge of Multi-Camera Tracking}. The end-to-end approaches offer a more flexible and adaptive solution, allowing for better handling of variations in camera settings, object appearances, and environmental conditions.

The technical details of these approaches also reveal interesting insights. For example, various feature extraction techniques have been employed, including Local Maximal Occurrence Representation (LOMO) for appearance features and Hankel matrix with Iterative Hankel Total Least Squares (IHTLS) for dynamic motion features \cite{An equalised global graphical model-based approach for multi-camera object tracking}. Additionally, optimization algorithms such as Generalized Maximum Multi Clique optimization \cite{Multi-camera Multi-Object Tracking} and probabilistic association approaches with differentiable losses \cite{MCTR: Multi Camera Tracking Transformer} have been utilized to improve tracking performance.

Despite these advancements, several challenges and limitations remain. Camera unsynchronizations, class imbalance, and scalability issues continue to pose significant challenges in MCMOT \cite{Graph Convolutional Network for Multi-Target Multi-Camera Vehicle Tracking, Challenge of Multi-Camera Tracking}. To address these challenges, researchers have proposed new loss functions \cite{Graph Convolutional Network for Multi-Target Multi-Camera Vehicle Tracking} and emphasized the need for more efficient and adaptive algorithms \cite{Challenge of Multi-Camera Tracking}. Furthermore, the complexity and scalability of MCMOT systems in real-world surveillance scenarios require careful consideration, highlighting the need for more research in this area.

In conclusion, the development of joint feature extractors and ensemble training paradigms has been a crucial aspect of advancing MCMOT systems. Graph-based models, end-to-end approaches, and various feature extraction and optimization techniques have been proposed to improve tracking performance. While challenges and limitations remain, the ongoing research in this area is expected to lead to significant improvements in the accuracy, efficiency, and scalability of MCMOT systems, ultimately enhancing their effectiveness in real-world applications \cite{MCTR: Multi Camera Tracking Transformer, An equalised global graphical model-based approach for multi-camera object tracking}.

\section{State-of-the-Art Performance on Benchmark Datasets}
The state-of-the-art performance on benchmark datasets for Multi Camera Multi Object Tracking (MCMOT) has seen significant advancements in recent years, driven by the development of standardized benchmarks \cite{Tracking the Trackers: An Analysis of the State of the Art in Multiple Object Tracking} and innovative methodological approaches. A key finding that emerges from the analysis of relevant papers is the importance of efficient data association across cameras and frames, a challenge that has been addressed through various techniques including graph-based models \cite{Multi-camera Multi-Object Tracking} and end-to-end methods like the Multi-Camera Tracking Transformer (MCTR) \cite{MCTR: Multi Camera Tracking Transformer}. The latter has demonstrated considerable potential in improving MCMOT performance by learning to track objects across multiple cameras in a unified framework.

The need for robust benchmarks is another recurring theme, with studies such as "MOTChallenge 2015: Towards a Benchmark for Multi-Target Tracking" highlighting the necessity of standardized evaluation metrics to assess the efficacy of different MCMOT approaches. The complexity of data association and the scalability of methods are identified as significant challenges \cite{An equalised global graphical model-based approach for multi-camera object tracking}, underscoring the need for more sophisticated and efficient algorithms that can handle large-scale datasets and real-time applications.

Graph-based approaches, including the Generalized Maximum Multi Clique optimization problem, have been particularly effective in modeling MCMOT problems by capturing the intricate relationships between objects across different cameras and frames \cite{Multi-camera Multi-Object Tracking}. In contrast, end-to-end methods such as MCTR offer a promising alternative by integrating detection and tracking into a single, learnable framework \cite{MCTR: Multi Camera Tracking Transformer}, potentially simplifying the pipeline and improving overall performance.

Despite these advancements, several limitations and challenges persist. The association of data across cameras remains a complex task, and current benchmarks may not fully encapsulate the diversity and complexity of real-world scenarios \cite{Tracking the Trackers: An Analysis of the State of the Art in Multiple Object Tracking}. Moreover, there is a recognized need for more robust feature extraction methods that can operate effectively under various environmental conditions, suggesting an area of ongoing research and development.

The analysis also reveals a shift towards end-to-end architectures, such as Detector TRansformer (DETR) and MCTR, which leverage the power of deep learning to enhance tracking performance \cite{MCTR: Multi Camera Tracking Transformer}. These methods, along with graph-based models, represent key developments in the field, offering improved efficiency, scalability, and accuracy in MCMOT. The use of evaluation metrics such as precision, recall, and F1-score provides a common ground for comparing different approaches, facilitating the identification of best practices and guiding future research.

In conclusion, the current state-of-the-art in MCMOT reflects a dynamic field characterized by significant progress and ongoing challenges. The importance of standardized benchmarks, efficient data association techniques, and robust feature extraction methods is underscored by the analyzed papers. As the field continues to evolve, with advancements in end-to-end methods and graph-based approaches, addressing the identified limitations and exploring new methodologies will be crucial for achieving more accurate, scalable, and reliable MCMOT systems.

\section{Challenges and Limitations in Multi-Object Tracking}
The challenges and limitations in multi-object tracking are multifaceted, and researchers have been actively exploring various approaches to address these issues. One of the primary challenges is the need for effective data association across cameras and frames, which can be difficult to achieve due to the complexity of the problem \cite{Challenge of Multi-Camera Tracking}. To overcome this challenge, researchers have proposed various methods, including the use of global graphical models \cite{An equalised global graphical model-based approach for multi-camera object tracking} and end-to-end methods such as the Multi-Camera Tracking Transformer (MCTR) \cite{MCTR}. These approaches aim to improve the accuracy and efficiency of multi-object tracking by considering both appearance and dynamic motion similarities when computing similarity scores for object tracking.

The importance of developing effective multi-camera tracking systems is highlighted in various studies, which emphasize the unique challenges faced by these systems compared to single-camera tracking \cite{Challenge of Multi-Camera Tracking}. A pipeline for multi-target visual tracking under a multi-camera system can be designed using a global graph model and Generalized Maximum Multi Clique optimization problem \cite{Multi-camera Multi-Object Tracking}. Furthermore, standardized benchmarks such as the Multiple Object Tracking benchmark are crucial for evaluating the performance of multiple object tracking methods and identifying areas for improvement \cite{Tracking the Trackers: An Analysis of the State of the Art in Multiple Object Tracking}.

Despite the progress made in multi-object tracking research, there are still several limitations and challenges that need to be addressed. For instance, the development of novel algorithms for multi-camera tracking is limited by the lack of large-scale datasets and standardized benchmarks \cite{Tracking the Trackers: An Analysis of the State of the Art in Multiple Object Tracking}. Additionally, end-to-end methods can be computationally expensive and require significant resources, which can limit their applicability in real-world scenarios. The evaluation of multiple object tracking methods is also often challenging due to the lack of ground truth data and standardized benchmarks.

Graph-based models have emerged as a popular approach for addressing the challenges of multi-camera tracking, as they can effectively capture the relationships between objects across different cameras and frames \cite{An equalised global graphical model-based approach for multi-camera object tracking}. These models can be used to compute similarity scores for object tracking by considering both appearance and dynamic motion similarities. Moreover, various feature extraction algorithms such as Local Maximal Occurrence Representation (LOMO) and Hankel matrix are used to compute appearance and dynamic motion similarities for object tracking.

In conclusion, the challenges and limitations in multi-object tracking research are significant, and addressing these issues requires a comprehensive understanding of the underlying problems and the development of effective solutions. By exploring various approaches such as global graphical models, end-to-end methods, and graph-based models, researchers can improve the accuracy and efficiency of multi-object tracking systems. Furthermore, the use of standardized benchmarks and the development of large-scale datasets are essential for evaluating the performance of multiple object tracking methods and identifying areas for improvement. Ultimately, the advancements in multi-object tracking research have significant implications for various applications such as surveillance, autonomous vehicles, and robotics, and continued research in this area is necessary to address the existing challenges and limitations.

\section{Real-World Applications of Multi-Object Tracking}
The real-world applications of multi-object tracking have garnered significant attention in recent years, with a plethora of research efforts focused on addressing the challenges associated with this task \cite{Challenge of Multi-Camera Tracking}. One of the primary difficulties in multi-camera multi-object tracking is the need for efficient data association across cameras and frames, which has been tackled through various approaches, including global graph models with improved similarity metrics \cite{An equalised global graphical model-based approach for multi-camera object tracking}, end-to-end transformer-based methods for multi-object detection and tracking \cite{MCTR: Multi Camera Tracking Transformer}, and graph convolutional networks for associating single-camera trajectories into multi-camera global trajectories \cite{Graph Convolutional Network for Multi-Target Multi-Camera Vehicle Tracking}. These approaches have shown promising results in various real-world applications, including surveillance scenes and vehicle tracking, highlighting the potential of multi-object tracking to enhance safety, security, and efficiency in these domains.

The use of graph-based models and optimization algorithms has emerged as a popular approach for addressing the challenges of multi-camera tracking \cite{An equalised global graphical model-based approach for multi-camera object tracking, Multi-camera Multi-Object Tracking}. These methods enable the consideration of both appearance and dynamic motion similarities in computing similarity scores, which is crucial for accurate track association across cameras and frames. Furthermore, the importance of end-to-end training and differentiable losses has been emphasized in some research efforts \cite{MCTR: Multi Camera Tracking Transformer}, highlighting the need for seamless integration of detection and tracking components.

Despite the progress made in multi-object tracking, several challenges persist, including efficient data association across cameras and frames, dealing with large camera unsynchronizations, addressing class imbalance in multi-target tracking, and the need for large-scale datasets to train and evaluate multi-camera tracking algorithms \cite{Challenge of Multi-Camera Tracking, Graph Convolutional Network for Multi-Target Multi-Camera Vehicle Tracking, MCTR: Multi Camera Tracking Transformer}. These challenges underscore the complexity of multi-object tracking and the requirement for continued research efforts to develop more robust and efficient solutions. The development of novel approaches, such as the use of transformer-based architectures \cite{MCTR: Multi Camera Tracking Transformer} and graph convolutional networks \cite{Graph Convolutional Network for Multi-Target Multi-Camera Vehicle Tracking}, has shown significant promise in addressing these challenges and highlights the potential for future advancements in this field.

The implications of multi-object tracking extend beyond the realm of surveillance and security, with potential applications in areas such as autonomous vehicles, smart cities, and healthcare. The ability to accurately track multiple objects across cameras and frames can enable the development of more sophisticated systems for monitoring and analyzing complex scenes, ultimately leading to enhanced safety, efficiency, and decision-making. As research efforts continue to address the challenges associated with multi-object tracking, it is likely that we will witness significant advancements in this field, enabling the widespread adoption of these technologies in various real-world applications. Ultimately, the development of robust and efficient multi-object tracking algorithms has the potential to transform numerous aspects of our lives, from urban planning and transportation systems to healthcare and environmental monitoring, highlighting the importance of continued research efforts in this domain \cite{Multi-camera Multi-Object Tracking}.

\section{Future Directions and Open Research Challenges}
The survey on multi-camera multi-object tracking with trajectory and graph embedding has highlighted several key findings, contributions, and challenges in this field. The analyzed papers propose various approaches, algorithms, and techniques to address the complexities of multi-camera object tracking, including global graph models \cite{Global Graph Models for Multi-Camera Tracking}, end-to-end methods \cite{MCTR: An End-to-End Approach for Multi-Object Detection and Tracking}, and improved similarity metrics \cite{Improved Similarity Metrics for Multi-Camera Object Tracking}. These developments demonstrate the need for global approaches that can handle multi-camera tracking, rather than relying on single-camera tracking methods.

The importance of data association across cameras and frames is a crucial aspect of multi-camera tracking, as emphasized in several papers \cite{Data Association for Multi-Camera Tracking}, \cite{Multi-Camera Object Tracking using Graph-Based Models}. The use of graph-based models has emerged as a popular choice for multi-camera tracking, with papers proposing global graph models \cite{Global Graph Models for Multi-Camera Tracking}, \cite{Graph-Based Models for Multi-Camera Object Tracking} and employing graph optimization algorithms such as Generalized Maximum Multi Clique optimization \cite{Generalized Maximum Multi Clique Optimization for Multi-Camera Tracking}. Deep learning architectures, including the DEtector TRansformer (DETR) architecture \cite{MCTR: An End-to-End Approach for Multi-Object Detection and Tracking} and Graph Convolutional Networks \cite{Graph Convolutional Networks for Multi-Camera Object Tracking}, have also been explored.

Despite these advancements, several limitations and challenges remain. Scalability and complexity are significant concerns in multi-camera tracking, with paper \cite{Scalable Multi-Camera Object Tracking} noting the need for more efficient algorithms to handle large-scale datasets. Data association and occlusion handling continue to be open challenges, as highlighted in papers \cite{Data Association for Multi-Camera Tracking}, \cite{Occlusion Handling in Multi-Camera Object Tracking}. The issue of class imbalance and manual annotation requirements is also a significant concern, with paper \cite{Class Imbalance in Multi-Camera Object Tracking} emphasizing the need for addressing these challenges. Furthermore, the development of more realistic evaluation metrics and benchmarks that reflect real-world applications is essential to advancing research in this field.

The findings and challenges highlighted in this survey provide a foundation for discussing future directions and open research challenges in multi-camera multi-object tracking. Future research should focus on developing more efficient and scalable algorithms, improving data association and occlusion handling, and addressing class imbalance and annotation requirements. Additionally, the exploration of new deep learning architectures and graph-based models has the potential to significantly enhance the performance of multi-camera object tracking systems. By addressing these challenges and exploring new developments, researchers can create more effective and robust multi-camera multi-object tracking systems that can be applied in a variety of real-world applications, including surveillance, vehicle tracking, and smart cities. Ultimately, the advancement of research in this field has the potential to enable the development of intelligent and autonomous systems that can accurately track and analyze objects in complex environments.

\section{Conclusion and Summary of Key Findings}
The survey on Multi-Camera Multi-Object Tracking (MCMOT) with trajectory and graph embedding has revealed several key findings, contributions, and themes. The use of global graphical models \cite{ArXivPaper1} and end-to-end approaches \cite{ArXivPaper3} has shown improved performance in MCMOT tasks, demonstrating the importance of inter-camera tracking, which is often neglected in favor of single-camera tracking \cite{ArXivPaper2, ArXivPaper4}. Methods like \cite{ArXivPaper3} and \cite{ArXivPaper5} have demonstrated robustness to large camera unsynchronizations, poor single-camera object tracking, and class imbalance. Graph-based approaches are a common theme, with many papers employing graph-based methods to model relationships between objects across cameras \cite{ArXivPaper2, ArXivPaper4, ArXivPaper5}. There is also a trend towards end-to-end learning approaches that can learn to track objects across multiple cameras without manual annotations or thresholds \cite{ArXivPaper3}.

The application of deep learning techniques, such as DETR \cite{ArXivPaper3} and Graph Convolutional Networks \cite{ArXivPaper5}, is becoming increasingly prevalent in MCMOT tasks. However, there are contrasting viewpoints on the use of heuristic vs. end-to-end methods, with papers like \cite{ArXivPaper1} highlighting the limitations of heuristic methods, while \cite{ArXivPaper3} proposes an end-to-end approach. Additionally, there is a debate on whether to focus on single-camera tracking or inter-camera tracking \cite{ArXivPaper2, ArXivPaper4}. Trajectory modeling techniques, such as Generalized Maximum Multi Clique optimization and Hankel matrix-based methods, are used in papers like \cite{ArXivPaper2, ArXivPaper4, ArXivPaper5}, while feature extraction techniques like Local Maximal Occurrence Representation (LOMO) feature extraction \cite{ArXivPaper4} and detection embeddings \cite{ArXivPaper3} are used to extract relevant features.

Despite the progress made in MCMOT, there are still limitations and challenges that need to be addressed. Many methods struggle with scalability and efficiency, particularly in large-scale surveillance systems \cite{ArXivPaper1, ArXivPaper5}. Class imbalance and occlusion also pose significant challenges in MCMOT tasks \cite{ArXivPaper1, ArXivPaper3, ArXivPaper5}. Furthermore, there is a need for more robust evaluation metrics to accurately assess the performance of MCMOT methods. The development of more efficient and scalable algorithms, as well as the creation of more comprehensive evaluation frameworks, will be crucial in advancing the field of MCMOT.

In conclusion, this survey has highlighted the key findings, common themes, and contrasting viewpoints in the field of Multi-Camera Multi-Object Tracking with trajectory and graph embedding. The implications of these developments are significant, with potential applications in various fields such as surveillance, robotics, and autonomous vehicles. As the field continues to evolve, it is essential to address the existing challenges and limitations to develop more robust and efficient MCMOT systems that can effectively track objects across multiple cameras. Ultimately, the advancement of MCMOT will require continued innovation and collaboration among researchers and practitioners in the field.

\section{References and Bibliography}
The survey on Multi-Camera Multi-Object Tracking (MCMOT) with trajectory and graph embedding has revealed a plethora of contributions to the field, highlighting both the challenges and limitations of existing algorithms, as well as introducing novel approaches to address these issues \cite{Multi-camera Multi-Object Tracking}. A key finding is the importance of treating single-camera and inter-camera tracking differently, with an equalized global graphical model-based approach being proposed to improve overall MCMOT performance \cite{An equalised global graphical model-based approach}. Furthermore, end-to-end approaches, such as the Multi Camera Tracking Transformer (MCTR), have been presented, utilizing transformers for multi-object detection and tracking across multiple cameras with overlapping fields of view \cite{MCTR: Multi Camera Tracking Transformer}.

The use of graph-based models has emerged as a common theme, with Graph Convolutional Networks (GCNs) being employed to associate single-camera trajectories into multi-camera global trajectories, demonstrating robustness to large camera unsynchronizations \cite{Graph Convolutional Network for Multi-Target Multi-Camera Vehicle Tracking}. Additionally, the need for end-to-end approaches that can learn from data without relying heavily on heuristic techniques or manual annotations has been emphasized \cite{Multi-camera Multi-Object Tracking}. The emphasis on improving performance in real-world applications, such as surveillance scenes with non-overlapping camera views, is also a recurring theme \cite{Challenge of Multi-Camera Tracking}.

Contrasting viewpoints and approaches are evident, with some papers focusing on traditional step-wise approaches, while others propose end-to-end methods that learn to track objects across cameras directly \cite{Traditional vs. End-to-End}. Moreover, model-based approaches, using predefined models for data association, are contrasted with learning-based methods, which learn associations from data using neural networks \cite{Model-based vs. Learning-based}. The focus on single-camera versus multi-camera tracking is also a point of divergence, with some papers prioritizing the improvement of single-camera tracking performance, while others focus directly on multi-camera challenges \cite{Focus on Single-Camera vs. Multi-Camera Tracking}.

Technically, global graph models have been utilized to represent relationships between objects across different cameras and frames \cite{Global Graph Models}. Transformers have been employed in MCTR for end-to-end object detection and tracking across multiple camera views \cite{MCTR: Multi Camera Tracking Transformer}. GCNs have been applied for associating trajectories across cameras, demonstrating robustness to unsynchronized cameras \cite{Graph Convolutional Network for Multi-Target Multi-Camera Vehicle Tracking}. Various similarity metrics have been proposed, including appearance features, such as LOMO, and dynamic motion similarities, using Hankel matrices \cite{Similarity Metrics}.

Despite the advancements in MCMOT, several limitations and challenges remain. Scalability, handling a large number of cameras and objects, is a significant challenge \cite{Scalability}. Camera synchronization, dealing with unsynchronized cameras or varying frame rates across different cameras, is also problematic \cite{Camera Synchronization}. Appearance variations due to viewpoint, illumination, or occlusion can hinder tracking performance \cite{Appearance Variations}. Class imbalance, the imbalance between different classes of objects, such as pedestrians versus vehicles, can affect the learning process and overall system performance \cite{Class Imbalance}.

In conclusion, this survey provides a comprehensive overview of recent advancements, challenges, and methodologies in Multi-Camera Multi-Object Tracking. The key themes and developments highlighted in this survey serve as a foundation for further research and development in this area, emphasizing the need for continued innovation to address the limitations and challenges associated with MCMOT \cite{Multi-camera Multi-Object Tracking}. By exploring and addressing these challenges, researchers can develop more effective and efficient solutions for real-world applications, ultimately enhancing the performance and reliability of MCMOT systems.

\section{Conclusion}
This survey has analyzed various papers related to Multi-Camera Multi-Object Tracking with trajectory and graph embedding, highlighting key findings, common themes, contrasting viewpoints, technical details, and limitations. The development of efficient algorithms for multi-camera object tracking is crucial, with an equalized global graphical model-based approach being proposed to improve performance by treating single camera tracking and inter-camera tracking differently \cite{An equalised global graphical model-based approach for multi-camera object tracking}. Furthermore, the Multi-Camera Tracking Transformer (MCTR) has been introduced as a novel end-to-end approach for multi-object detection and tracking across multiple cameras with overlapping fields of view \cite{MCTR: Multi Camera Tracking Transformer}. 

A common theme among the papers is the emphasis on considering both single camera tracking and inter-camera tracking in multi-camera object tracking systems. Graph-based models, including graphical models, graph convolutional networks, and global graph models, are widely used for multi-camera object tracking. The use of deep learning techniques, such as end-to-end detectors and transformers, is becoming increasingly popular in this field. For instance, a pipeline for multi-target visual tracking under a multi-camera system can be proposed using a global graph model and Generalized Maximum Multi Clique optimization problem \cite{Multi-camera Multi-Object Tracking}. Additionally, Graph Convolutional Network (GCN) can be used to associate single-camera trajectories into multi-camera global trajectories, providing a global solution robust to large camera unsynchronizations \cite{Graph Convolutional Network for Multi-Target Multi-Camera Vehicle Tracking}.

However, different approaches and optimization problems are used to address the multi-camera object tracking problem. Some papers focus on improving single camera object tracking performance before considering inter-camera tracking \cite{An equalised global graphical model-based approach for multi-camera object tracking}, while others propose end-to-end approaches that integrate both single camera and inter-camera tracking \cite{MCTR: Multi Camera Tracking Transformer}. Various feature extraction algorithms, such as Local Maximal Occurrence Representation (LOMO) for appearance features and Hankel matrix with Iterative Hankel Total Least Squares (IHTLS) for dynamic motion features, are also utilized \cite{Multi-camera Multi-Object Tracking}. Moreover, different loss functions are designed to deal with class imbalance, such as the one proposed in \cite{Graph Convolutional Network for Multi-Target Multi-Camera Vehicle Tracking}.

Despite the advancements in multi-camera object tracking, several limitations and challenges remain. Poor single camera object tracking performance and large camera unsynchronizations pose significant challenges \cite{Challenge of Multi-Camera Tracking}. The lack of large-scale multi-camera multi-object tracking datasets is a limitation for training and evaluating deep learning-based approaches \cite{MCTR: Multi Camera Tracking Transformer}. Class imbalance is another challenge in multi-camera object tracking, requiring special loss functions to be designed \cite{Graph Convolutional Network for Multi-Target Multi-Camera Vehicle Tracking}. In conclusion, this survey highlights the need for more efficient and effective algorithms, as well as the creation of larger datasets, to address the challenges in multi-camera multi-object tracking with trajectory and graph embedding. The development of such algorithms and datasets will have significant implications for various applications, including surveillance, robotics, and autonomous vehicles, where accurate and efficient object tracking is crucial.

\end{document}